\cleardoublepage
\chapter*{历史注记}
\label{nota-historica}
\addcontentsline{toc}{chapter}{历史注记}
\markboth{历史注记}{}

\thispagestyle{empty}

\noindent 本书是一部虚构作品。若阿金·安东尼奥·科雷亚、患者1,260和他携带的怪物从未存在。被叙述为\textit{收容院}的机构是虚构的。

\vspace{0.8em}

\noindent 但这一切虚构所依托的地基是真实的——正因如此，它必须在这里结束，而不是以一次惊吓结束。

\vspace{0.8em}

\noindent 1915年大旱期间，塞阿拉政府在福塔莱萨附近建立了\textbf{阿拉加迪索（Alagadiço）收容营}，逃荒者在那里遭到围禁。塞纳多尔·蓬佩乌的\textbf{帕图（Patu）收容营}属于另一个历史阶段：1932年大旱，其运作时间为1932--1933年。本版以阿拉加迪索作为1915年场景的历史参照；若阿金、他的家人、乱葬坑以及1917年的转送均为虚构元素。收容与死亡人数不作为精确数字呈现。

\vspace{0.8em}

\noindent 米纳斯吉拉斯州的\textbf{巴巴塞纳收容院}真实存在。二十世纪期间，估计有\textbf{六万人}在其墙内死亡——其中许多人根本没有精神科诊断：他们是穷人、黑人、癫痫患者、单亲母亲、同性恋者、“碍事”的人。有无麻醉的电休克、束缚、饥饿、乱葬坑，以及向医学院出售遗体。这些都不是本书杜撰的；它们是有据可查的历史——参见达妮埃拉·阿尔贝克斯的\textit{巴西大屠杀}等资料。

\vspace{0.8em}

\noindent 这些书页的超自然恐怖是一种隐喻。那种跨越世代、无人处理、像传染病一样传递的痛苦，才是本书唯一的真实怪物——它有名字：遗忘。对牲畜圈和收容院死者最糟糕的做法，不是用一个故事去闹鬼。而是不给他们留下任何故事。

\vspace{0.8em}

\noindent \textbf{关于审美化的风险:} 作者意识到，创伤文学行走在象征性加工与他人苦难消费之间的边界上。本书无意将收容院的恐怖作为奇观呈现，也无意取代真实受害者及其后代的证词——这些证词不可替代，且优先于任何虚构。如果阅读这些书页只产生惊吓，它就失败了；如果它在面对遗忘时产生不适，它就完成了自己的功能。对感到受影响的读者：中断阅读永远是正当的，支持渠道可用——CVV 188（巴西，24小时）、988 自杀与危机生命线（美国），以及中国心理援助热线 12356。

\vspace{1.2em}

\noindent\hfill\textit{愿时间埋葬的，不被抹去。}\hspace{2em}
